\section{Limitations, conclusions and unresolved questions}
\label{sec:limits}

\paragraph{Limitations.}
\begin{itemize}\itemsep0pt
\item \emph{No retraining.} I tested the paper's claims through released checkpoints. The checkpoints differ in more than the one factor I compare (fusion, AudioSet+K2C), so E2 and E3 are observational, not controlled ablations.
\item \emph{One dataset, synthetic noise.} All results are on ESC-50 with Gaussian coloured noise. Real background noise (speech babble, traffic, music) is non-stationary and was not tested. The brown-noise condition turned out to be almost clean in the model's frequency band, so only two noise types are informative.
\item \emph{SNR definition.} SNR is computed over the whole clip; for clips with short events the local SNR during the event is higher than nominal.
\item \emph{Possible train/test overlap.} ESC-50 and LAION-Audio-630K both draw on Freesound. The authors state that overlap was handled, but I could not verify it for the released checkpoints, so absolute zero-shot numbers may be optimistic. This does not affect the \emph{paired} comparisons between adapters, which share the same embeddings.
\item \emph{Uncertainty.} Intervals are over 15 (fold, seed) runs that are not independent; they describe run-to-run variation, not population uncertainty.
\item \emph{Hyper-parameter grid.} The pre-specified $\alpha$ grid was too narrow. The follow-up widened it after seeing the boundary (selection still on validation data only), but even the wider grid was often at its edge.
\item \emph{Only one backbone and one adapter family.} Treff's learned cross-attention, Tip-Adapter-F and noise-augmented fine-tuning of the probe or encoder were not tried.
\end{itemize}

\paragraph{Conclusions.}
(1) The released LAION-CLAP checkpoints reproduce, and slightly exceed, the paper's ESC-50 zero-shot accuracy (\repaudiosetfusionA--\repaudiosetA\% vs.\ 88.0--91.0\%), but the result moves by several points with the prompt and by about two points with the evaluation pipeline.
(2) Feature fusion cannot matter on ESC-50 because the code bypasses it for clips of 10\,s or less; the measured differences between fusion and non-fusion checkpoints are within fold noise.
(3) Zero-shot CLAP loses about 25 points at 0\,dB white or pink noise and collapses onto noise-like classes.
(4) Adding noise-augmented copies of the support clips to a Tip-Adapter cache improves noisy accuracy over a clean cache by about 2--3.5 points at $K{=}16$, also for a noise colour that was not used to build the cache, with no measurable clean-accuracy cost and negligible run-time cost.
(5) The improvement comes from the augmented data rather than from the cache: a logistic-regression probe trained on the same data gains more and is the best method for $K\ge4$. The training-free cache is preferable only with one clip per class.

\paragraph{Unresolved questions.}
Why does the validation fold keep choosing a cache weight that overrides the text prior, and would a calibrated combination of probe and zero-shot logits beat both?
Does the gain hold for real recorded noise and for longer audio where fusion is actually used?
Can the augmentation SNR range be chosen without knowing the deployment noise, for example by sampling several noise colours?
